\documentclass[10pt,a4paper]{amsart}
\usepackage{amsmath,amssymb,amsthm,graphicx}
\usepackage{amsmath}
\usepackage{amssymb}
\usepackage{hyperref}
% Class-compatibility preamble for a reader excerpt of
%   Paul Balmer, "Homological support of big objects in tensor-triangulated categories",
%   Journal de l'Ecole polytechnique -- Mathematiques 7 (2020), DOI 10.5802/jep.135.
%
% This file is the source directory's inherited layout adapter, unchanged except for one
% mechanical line: \smaller is declared with \providecommand instead of \newcommand, so that
% the excerpt compiles whether or not the wrapper class already defines \smaller as a
% font-size command. No mathematical macro, symbol, statement, proof step or reference was
% added, changed or removed by this file.
\usepackage{amsmath,amssymb,amsthm,xfrac,xspace,hyperref,graphicx,comment}
\excludecomment{DefTralics}
\providecommand{\C}{}
\providecommand{\bibliofont}{\small}
\newcommand{\ie}{i.e.,\xspace}
\newcommand{\eg}{e.g.,\xspace}
\newcommand{\etc}{etc.\xspace}
\newcommand{\cf}{cf.\xspace}
\newcommand{\resp}{resp.\xspace}
\providecommand{\smaller}{\small}
\newcommand{\Subsection}{\subsection}
\newcommand{\Subsubsection}{\subsubsection}
\newcommand{\skpt}{\smallskip}
\emergencystretch=3em
\providecommand{\loccit}{\textit{loc.\ cit.}}
\newenvironment{enumeratea}
{\bgroup\def\theenumi{\alph{enumi}}\begin{enumerate}}
{\end{enumerate}\egroup}
\newenvironment{enumeratei}
{\bgroup\def\theenumi{\roman{enumi}}\begin{enumerate}}
{\end{enumerate}\egroup}

\def\hto{\mathrel{\lhook\joinrel\to}}
\def\mto{\mathchoice{\longmapsto}{\mapsto}{\mapsto}{\mapsto}}
\def\To#1{\mathchoice{\xrightarrow{\textstyle\kern4pt#1\kern3pt}}{\stackrel{#1}{\longrightarrow}}{}{}}

\newcommand{\onto}{\mathchoice{\to\hspace*{-5mm}\to}{\to\hspace*{-2.85mm}\to}{\to\hspace*{-3mm}\to}{\to\hspace*{-3mm}\to}}

\newcommand{\into}{\mathchoice{\xymatrix@=.5cm{\ar@{>->}[r]&}}
{\xymatrix@=.3cm{\ar@{>->}[r]&}\hspace*{-1mm}}
{\hspace*{-1mm}\xymatrix@=.25cm{\ar@{>->}[r]&}\hspace*{-1.5mm}}
{\xymatrix@=.1cm{\ar@{>->}[r]&}\hspace*{-1mm}}
}
\usepackage[mathscr]{eucal}
\usepackage{bbm}
\usepackage{cleveref}
\crefdefaultlabelformat{#2\textup{#1}#3}
\usepackage[all,cmtip]{xy}\let\labelstyle\textstyle
\newdir{ >}{{}*!/-10pt/\dir{>}}
\newdir^{ (}{{}*!/-5pt/\dir^{(}}
\newdir_{ (}{{}*!/-5pt/\dir_{(}}

\numberwithin{equation}{section}

\swapnumbers

\newtheorem{Thm}[equation]{Theorem}
\newtheorem*{Thm*}{Theorem}
\newtheorem{Prop}[equation]{Proposition}
\newtheorem{Lem}[equation]{Lemma}
\newtheorem{Cor}[equation]{Corollary}

\theoremstyle{definition}
\newtheorem{Def}[equation]{Definition}
\newtheorem{Exa}[equation]{Example}
\newtheorem{Exas}[equation]{Examples}
\newtheorem{Cons}[equation]{Construction}
\newtheorem{Hyp}[equation]{Hypothesis}
\newtheorem{Rem}[equation]{Remark}
\newtheorem{Rec}[equation]{Recollection}

\DeclareMathOperator{\coker}{coker}
\DeclareMathOperator{\colim}{colim}
\DeclareMathOperator{\cone}{cone}
\DeclareMathOperator{\End}{End}
\DeclareMathOperator{\Hom}{Hom}
\DeclareMathOperator{\id}{id}
\DeclareMathOperator{\Id}{Id}
\DeclareMathOperator{\Ind}{Ind}
\DeclareMathOperator{\Ker}{Ker}
\DeclareMathOperator{\opname}{op}
\DeclareMathOperator{\supp}{supp}
\DeclareMathOperator{\Spc}{Spc}
\DeclareMathOperator{\Spec}{Spec}

\newcommand{\bbZ}{\mathbb{Z}}
\newcommand{\cat}[1]{\mathscr{#1}}
\newcommand{\Endcat}[1]{\End_{\cat #1}}
\newcommand{\Homcat}[1]{\Hom_{\cat #1}}
\newcommand{\inv}{^{-1}}
\newcommand{\isoto}{\overset{\sim}{\longrightarrow}}
\newcommand{\isotoo}{\overset{\sim}{\longrightarrow}}
\newcommand{\op}{^{\opname}}
\let\xto\To
\newcommand{\oto}[1]{\overset{#1} \longrightarrow}
\newcommand{\otoo}[1]{\overset{#1}{\longrightarrow}}
\newcommand{\Paul}[1]{{\color{ForestGreen}#1}}
\newcommand{\SET}[2]{\big\{#1\,\big|\,#2\big\}}
\newcommand{\too}{\mathop{\longrightarrow}\limits}
\newcommand{\unit}{\mathbbm{1}}

\DeclareMathOperator{\Ab}{Ab}
\DeclareMathOperator{\Add}{Add}
\DeclareMathOperator{\coev}{coev}
\DeclareMathOperator{\Der}{D}
\DeclareMathOperator{\Img}{Im}
\DeclareMathOperator{\Loc}{Loc}
\DeclareMathOperator{\modname}{mod}
\DeclareMathOperator{\Mod}{Mod}
\DeclareMathOperator{\SH}{SH}
\newcommand{\smallh}{\mathrm{h}}
\DeclareMathOperator{\Stab}{Stab}
\DeclareMathOperator{\Supp}{Supp}

\newcommand{\ababs}{\textit{ab absurdo}}
\newcommand{\ABp}{\bar{\cA}_{\cB'}}
\newcommand{\AB}{\bar{\cA}_{\cB}}
\newcommand{\AC}{\bar{\cA}_{\cC}}
\newcommand{\adj}{\dashv}
\newcommand{\bat}[1]{\bar{\cat{#1}}}
\newcommand{\bEB}{\bar E_{\cB}}
\newcommand{\bEC}{\bar E_{\cC}}
\newcommand{\boneda}{\bar\yoneda}
\newcommand{\bunit}{\bar\unit}
\newcommand{\cA}{\cat{A}}
\newcommand{\cB}{\cat{B}}
\newcommand{\cC}{\cat{C}}
\newcommand{\cD}{\cat{D}}
\newcommand{\cP}{\cat{P}}
\newcommand{\cS}{\cat{S}}
\newcommand{\cT}{\cat{T}}
\newcommand{\EBp}{E_{\cB'}}
\newcommand{\EB}{E_{\cB}}
\newcommand{\EC}{E_{\cC}}
\newcommand{\fp}{^{\textup{fp}}}
\newcommand{\hBp}{\boneda_{\cB'}}
\newcommand{\hB}{\boneda_{\cB}}
\newcommand{\hC}{\boneda_{\cC}}
\newcommand{\hEB}{\hat E_{\cB}}
\newcommand{\hEC}{\hat E_{\cC}}
\let\hook\hto
\newcommand{\hunit}{\hat\unit}
\newcommand{\ideal}[1]{\langle #1\rangle}
\newcommand{\modd}{\modname\text{\!-}}
\newcommand{\Modd}{\Mod\text{\!-}}
\newcommand{\MS}{\Modd\cat S^c}
\newcommand{\mS}{\modd\cat S^c}
\newcommand{\MT}{\Modd\cat T^c}
\newcommand{\mT}{\modd\cat T^c}
\newcommand{\potimes}[1]{^{\otimes #1}}
\newcommand{\qquadtext}[1]{\qquad\text{#1}\qquad}
\newcommand{\restr}[1]{_{|_{\scriptstyle #1}}}
\newcommand{\sbull}{{\scriptscriptstyle\bullet}}
\newcommand{\SpchS}{\Spch(\cat{S}^c)}
\newcommand{\SpchT}{\Spch(\cat{T}^c)}
\newcommand{\Spch}{\Spc^{\smallh}}
\newcommand{\SpcS}{\Spc(\cat{S}^c)}
\newcommand{\SpcT}{\Spc(\cat{T}^c)}
\newcommand{\tensids}{$\otimes$\nobreakdash-ideals}
\newcommand{\tensid}{$\otimes$\nobreakdash-ideal}
\newcommand{\tensnil}{$\otimes$\nobreakdash-nilpotent}
\newcommand{\vcorrect}[1]{{\vphantom{\vbox to #1em{}}}}
\newcommand{\yoneda}{\mathbbm{h}}
\crefname{originaldefinition}{original definition}{original definitions}
\Crefname{originaldefinition}{Original definition}{Original definitions}
\makeatletter
\newlabel{Cons:E_C}{{6}{source}{}{sourcepointer-1}{}}
\newlabel{Cons:E_C@cref}{{[originaldefinition][0][]6}{[1][0][]source}{}{}{}}
\makeatother
\crefname{originalproposition}{original proposition}{original propositions}
\Crefname{originalproposition}{Original proposition}{Original propositions}
\makeatletter
\newlabel{Prop:A/C}{{4}{source}{}{sourcepointer-2}{}}
\newlabel{Prop:A/C@cref}{{[originalproposition][0][]4}{[1][0][]source}{}{}{}}
\makeatother
\crefname{originalproposition}{original proposition}{original propositions}
\Crefname{originalproposition}{Original proposition}{Original propositions}
\makeatletter
\newlabel{Prop:ring-in-AB}{{5}{source}{}{sourcepointer-3}{}}
\newlabel{Prop:ring-in-AB@cref}{{[originalproposition][0][]5}{[1][0][]source}{}{}{}}
\makeatother
\crefname{originalproposition}{original proposition}{original propositions}
\Crefname{originalproposition}{Original proposition}{Original propositions}
\makeatletter
\newlabel{it:A/C-d}{{4}{source}{}{sourcepointer-4}{}}
\newlabel{it:A/C-d@cref}{{[originalproposition][0][]4}{[1][0][]source}{}{}{}}
\makeatother
\title{Nested localising tensor ideals give isomorphic weak rings of the associated pure-injectives}
\author{Paul Balmer}
\date{doi:10.5802/jep.135; CC BY4.0}
\begin{document}
\maketitle
\noindent\textit{Editorial note.} All delivered mathematical text is the author's own, copied from the pinned publisher source sources/jep-2020-135-homological-support/full-original-native.tex; no statement or proof step was rewritten, reordered or removed.\par\medskip
\noindent\textit{Editorial note.} Preamble provenance: the retained preamble is the source's own retained-preamble.tex, which carries the author's macro and theorem declaration block verbatim (its own header records which lines of the pinned file it copies); the wrapper class supplies the baseline packages that block assumes. Nothing was added, no macro was invented, and the excerpt serves only as preamble, never as a span coordinate.\par\medskip
\noindent\textit{Editorial note.} External references are navigation pointers only; every deliberately kept step of the author's argument stays inside the counted spans. The displayed numbers of those pointers are the author's own numbering under the author's \textbackslash{}newtheorem declarations (the source has no subsections, so numbering is section.index), computed mechanically by the lane's scanner.\par\medskip
\noindent\textit{Editorial note.} Four references to other labelled items occur inside the counted spans (the construction Cons:E\_C and the propositions Prop:ring-in-AB, Prop:A/C and its clause it:A/C-d); each is recorded as a navigation stub with a kind from the assembler's enumeration and the author's own number. The labels containing a slash are exactly the ones that the assembler's whitelist now admits.\par\medskip
\noindent\textit{Editorial note.} The counted spans contain no bibliography citation, so the bibliography staged with the source is attached unchanged and no per-card excerpt was needed.\par\medskip
\noindent\textit{Editorial note.} Assembly history: the first attempt (r001) failed with 'Environment Prop undefined' because the preamble carried only this source's retained-preamble.tex, which for this source is a bare layout adapter and contains no  ewtheorem declarations; r002 adds the author's own macro and theorem declaration block (lines 6-143 of the pinned file). The rule was refined accordingly: a retained-preamble only replaces the author block when it actually contains theorem declarations. No statement, span or stub changed between the revisions.\par\medskip
\noindent\textit{Allowed standard prerequisites (editorial).}\begin{itemize}
\item The weak ring attached to a localising tensor ideal. Same source, Construction labelled Cons:E\_C; invoked in the counted statement at line 507. Supplies the objects EB and EC whose isomorphism is asserted.
\item Units of the weak rings in the intermediate quotient are monomorphisms. Same source, Proposition labelled Prop:ring-in-AB; invoked inside the counted proof at line 524. Identifies the unit of the intermediate weak ring with a monomorphism, which lets the injective hull of the unit be compared.
\item Isomorphisms in the intermediate quotient that descend to the big category. Same source, Proposition labelled Prop:A/C, clause (it:A/C-d); invoked at the end of the counted proof at line 533. Lifts the isomorphism produced in the intermediate quotient back to the ambient tensor-triangulated category and makes it compatible with the units.
\end{itemize}\medskip
\ifdefined\hypertarget\hypertarget{sourcepointer-1}{}\fi
\noindent\textit{Source reference.} Original definition 6: sources/jep-2020-135-homological-support/full-original-native.tex, source lines 507--507. Navigation pointer for \textbackslash{}ref\{Cons:E\_C\} used inside the counted span (line 507); the argument that consumes it is retained in the span.\par\medskip
\ifdefined\hypertarget\hypertarget{sourcepointer-2}{}\fi
\noindent\textit{Source reference.} Original proposition 4: sources/jep-2020-135-homological-support/full-original-native.tex, source lines 533--533. Navigation pointer for \textbackslash{}ref\{Prop:A/C\} used inside the counted span (line 533); the argument that consumes it is retained in the span.\par\medskip
\ifdefined\hypertarget\hypertarget{sourcepointer-3}{}\fi
\noindent\textit{Source reference.} Original proposition 5: sources/jep-2020-135-homological-support/full-original-native.tex, source lines 524--524. Navigation pointer for \textbackslash{}ref\{Prop:ring-in-AB\} used inside the counted span (line 524); the argument that consumes it is retained in the span.\par\medskip
\ifdefined\hypertarget\hypertarget{sourcepointer-4}{}\fi
\noindent\textit{Source reference.} Original proposition 4: sources/jep-2020-135-homological-support/full-original-native.tex, source lines 533--533. Navigation pointer for \textbackslash{}ref\{it:A/C-d\} used inside the counted span (line 533); the argument that consumes it is retained in the span.\par\medskip
\begin{Prop}
\label{Prop:EB=EC}
Let $\cB\in\SpchT$ and $\cC\subset\cA$ be a proper localizing \tensid, containing~$\cB$ or equivalently containing~$\Loc(\cB)$. Let $\EB$ and $\EC$ be the corresponding weak rings of \Cref{Cons:E_C}. Then there is an isomorphism $\EB\simeq \EC$ in~$\cT$ compatible with the unit maps~$\unit\to\EB$ and~$\unit\to \EC$.
\end{Prop}

\begin{proof}
Since $\Loc(\cB)\subseteq\cC$, we can perform the Gabriel quotient $\cA/\cC$ in two steps:
\[
\kern4em\xymatrix{
\cat{A} \ar@<-.3em>@{->>}[d]_-{Q_{\cB}} \ar@/_2em/@<-1em>@{->>}[dd]_-{Q_{\cC}}
&&& \hEC
\\
\AB \ar@<-.3em>@{_{ (}->}[u]_-{R_\cB} \ar@<-.3em>@{->>}[d]_-{\bar Q}
&&& \tilde E \ar@{|->}[u]
\\
\AC \ar@<-.3em>@{_{ (}->}[u]_-{\bar R} \ar@/_2em/@<-1em>@{_{ (}->}[uu]_-{R_{\cC}}
&&& \bEC \ar@{|->}[u]
}
\]
for an intermediate Gabriel localization~$\bar Q\adj\bar R$. Let $\tilde E:=\bar R(\bEC)$ the injective object in~$\AB$ associated to the injective hull~$\bunit\into \bEC$ in~$\AC$. By construction of~$\EC\in\cT$, we have $\hEC\cong R_{\cC}(\bEC)\cong R_{\cB}(\tilde E)$ and therefore $\tilde E\cong Q_\cB(\hEC)=\hB(\EC)$. Since $\EC$ is a weak ring in~$\cT$, its image $\tilde E$ is a weak ring in~$\AB$, that is non-zero since $\bar Q(\tilde E)\cong\bEC\neq0$. By \Cref{Prop:ring-in-AB}, the unit $\hB(\eta_{\cC})\colon\bunit\into \tilde E$ is therefore a monomorphism in~$\AB$. On~the other hand, $\hB(\eta_\cB)\colon\bunit\into\bEB$ is by definition the injective hull of~$\bunit$ in~$\AB$. Hence there exists a commutative diagram in~$\AB$
\begin{equation}
\label{eq:aux-max-vs-max}
\vcenter{\xymatrix@C=2cm{\unit \ar@{ >->}[r]^-{\hB(\eta_\cB)} \ar@<-.4ex>@{ >->}[rd]_-{\hB(\eta_{\cC})}
& \bEB \ar@{ >->}[d]^{\exists\,\varphi}
\\
& \tilde E
}}
\end{equation}
where $\varphi$ is a split monomorphism. Applying the exact functor~$\bar Q$ to this diagram and using that~$\bar Q(\tilde E)=\bEC$ is the injective hull of~$\bunit$ in~$\AC$, we see that $\bar Q(\varphi)$ must be an isomorphism. In summary, $\tilde E\cong \bEB\oplus N$ in~$\AB$ for an object~$N$ such that $\bar Q(N)=0$. But since $\tilde E=\bar R(\bEC)$ and $\bar R$ is fully faithful, this forces the object~$N$ to be in the image of~$\bar R$ as well (it is `$\bar R$-local'). It follows that $N\cong \bar R\bar Q(N)=0$. Therefore $\varphi$ is already an isomorphism $\bEB\simeq \tilde E\simeq \bar R(\bEC)$ in~$\AB$. By \Cref{Prop:A/C}\,\eqref{it:A/C-d} (applied to $\bat{A}=\cA/\Loc(\cB)$) the isomorphism $\varphi$ in~\eqref{eq:aux-max-vs-max} comes from~$\cT$ and is compatible with the units there.
\end{proof}

\begin{thebibliography}{99}
\bibitem{Balmer05a} P. Balmer - “The spectrum of prime ideals in tensor triangulated categories” , J. reine angew. Math. 588 (2005), p. 149-168
\bibitem{Balmer18a} P. Balmer - “On the surjectivity of the map of spectra associated to a tensor-triangulated functor” , Bull. London Math. Soc. 50 (2018) no. 3, p. 487-495
\bibitem{Balmer19} P. Balmer - “A guide to tensor-triangular classification” , in Handbook of homotopy theory (H. Miller, ed.), Chapman and Hall/CRC, 2019, Available on the author’s web page
\bibitem{Balmer20a} P. Balmer - “Nilpotence theorems via homological residue fields” , Tunis. J. Math. 2 (2020) no. 2, p. 359-378
\bibitem{BalmerCameron20pp} P. Balmer \& J. Cameron - “Computing homological residue fields in algebra and topology” , 2020
\bibitem{BalmerDellAmbrogioSanders16} P. Balmer, I. Dell’Ambrogio \& B. Sanders - “Grothendieck-Neeman duality and the Wirthmüller isomorphism” , Compositio Math. 152 (2016) no. 8, p. 1740-1776
\bibitem{BalmerFavi11} P. Balmer \& G. Favi - “Generalized tensor idempotents and the telescope conjecture” , Proc. London Math. Soc. (3) 102 (2011) no. 6, p. 1161-1185
\bibitem{BensonIyengarKrause08} D. J. Benson, S. B. Iyengar \& H. Krause - “Local cohomology and support for triangulated categories” , Ann. Sci. École Norm. Sup. (4) 41 (2008) no. 4, p. 573-619
\bibitem{BensonIyengarKrause11a} D. J. Benson, S. B. Iyengar \& H. Krause - “Stratifying modular representations of finite groups” , Ann. of Math. (2) 174 (2011) no. 3, p. 1643-1684
\bibitem{BensonIyengarKrause11c} D. J. Benson, S. B. Iyengar \& H. Krause - “Stratifying triangulated categories” , J. Topology 4 (2011) no. 3, p. 641-666
\bibitem{BensonIyengarKrause12b} D. J. Benson, S. B. Iyengar \& H. Krause - “Colocalizing subcategories and cosupport” , J. reine angew. Math. 673 (2012), p. 161-207
\bibitem{BensonIyengarKrause12a} D. J. Benson, S. B. Iyengar \& H. Krause - Representations of finite groups: local cohomology and support , Oberwolfach Seminars , vol. 43 , Birkhäuser/Springer, Basel, 2012
\bibitem{BensonIyengarKrause13} D. J. Benson, S. B. Iyengar \& H. Krause - “Module categories for group algebras over commutative rings” , J. K-Theory 11 (2013) no. 2, p. 297-329, With an appendix by Greg Stevenson
\bibitem{BalmerKrauseStevenson19} P. Balmer, H. Krause \& G. Stevenson - “Tensor-triangular fields: ruminations” , Selecta Math. (N.S.) 25 (2019) no. 1, article ID 13, 36 pages
\bibitem{BalmerKrauseStevenson20} P. Balmer, H. Krause \& G. Stevenson - “The frame of smashing tensor-ideals” , Math. Proc. Cambridge Philos. Soc. 168 (2020) no. 2, p. 323-343
\bibitem{DwyerPalmieri08} W. G. Dwyer \& J. H. Palmieri - “The Bousfield lattice for truncated polynomial algebras” , Homology Homotopy Appl. 10 (2008) no. 1, p. 413-436
\bibitem{HoveyPalmieriStrickland97} M. Hovey, J. H. Palmieri \& N. P. Strickland - Axiomatic stable homotopy theory , Mem. Amer. Math. Soc. , vol. 128, no. 610 , American Mathematical Society, Providence, RI, 1997
\bibitem{HoveyStrickland99} M. Hovey \& N. P. Strickland - Morava $K$-theories and localisation , Mem. Amer. Math. Soc. , vol. 139, no. 666 , American Mathematical Society, Providence, RI, 1999
\bibitem{Krause00} H. Krause - “Smashing subcategories and the telescope conjecture—an algebraic approach” , Invent. Math. 139 (2000) no. 1, p. 99-133
\bibitem{LurieHA} J. Lurie - “Higher algebra” (2017), Online at http://www.math.ias.edu/\textasciitilde{}lurie/
\bibitem{Neeman96} A. Neeman - “The Grothendieck duality theorem via Bousfield’s techniques and Brown representability” , J. Amer. Math. Soc. 9 (1996) no. 1, p. 205-236
\bibitem{Neeman00} A. Neeman - “Oddball Bousfield classes” , Topology 39 (2000) no. 5, p. 931-935
\bibitem{Neeman01} A. Neeman - Triangulated categories , Annals of Math. Studies , vol. 148 , Princeton University Press, Princeton, NJ, 2001
\bibitem{Stevenson13} G. Stevenson - “Support theory via actions of tensor triangulated categories” , J. reine angew. Math. 681 (2013), p. 219-254
\end{thebibliography}
\end{document}
